\documentclass[11pt,a4paper]{article}

\usepackage[T1]{fontenc}
\usepackage[utf8]{inputenc}
\usepackage[french]{babel}
\usepackage{lmodern}
\usepackage{microtype}
\usepackage{geometry}
\geometry{margin=2.5cm}
\usepackage{amsmath}
\usepackage{booktabs}
\usepackage{enumitem}
\usepackage{csquotes}
\usepackage{hyperref}
\hypersetup{
  colorlinks=true,
  linkcolor=black,
  citecolor=black,
  urlcolor=blue,
  pdftitle={Vers une sociologie internationale des personnes collectives},
  pdfauthor={Sylvain Geffroy}
}

\usepackage[backend=biber,style=authoryear,sorting=nyt,maxcitenames=2]{biblatex}
\addbibresource{../references/sipc_bibliographie.bib}

\title{\textbf{Vers une sociologie internationale des personnes collectives}\\[0.4em]
\large Champs, capitaux, dépendances, rangs et mécanismes de la société mondiale}
\author{Sylvain Geffroy}
\date{Version 1.0 --- \today}

\begin{document}
\maketitle

\begin{abstract}
\noindent
Cet article propose la \emph{Sociologie internationale des personnes collectives} (SIPC), un
cadre théorique et opérationnel qui analyse les relations internationales comme une société
mondiale stratifiée d'acteurs collectifs inégaux. La SIPC ne réduit ni la politique mondiale à
un système d'États rationnels dans l'anarchie, ni les acteurs à des récits désincarnés~: elle
les traite comme des \emph{personnes collectives} dotées de capitaux convertibles, de rôles
revendiqués et reconnus, de dépendances activables et de trajectoires dépendantes du chemin.
Son unité d'analyse n'est ni l'État ni la crise, mais la \emph{situation sociale
internationale}. Sa contribution propre n'est pas l'invention de concepts isolés --- la plupart
sont hérités de la sociologie des champs, de l'École anglaise, du constructivisme, de
l'économie politique internationale et de la sociologie analytique --- mais leur
\emph{intégration disciplinée} et leur \emph{formalisation en une ontologie testable}. Nous
exposons l'appareil conceptuel, sa logique causale configurationnelle et sa discipline
probatoire, puis l'appliquons à cinq situations contrastées (détroit de Taïwan, guerre en
Ukraine, rivalité technologique États-Unis / Chine, mer de Chine méridionale, Arctique). Nous
discutons enfin ses limites, notamment les risques d'anthropomorphisme, d'hypertrophie
conceptuelle et de fausse précision.

\medskip
\noindent\textbf{Mots-clés~:} relations internationales, sociologie, champs, capitaux,
dépendances, reconnaissance, mécanismes causaux, ontologie.
\end{abstract}

\section{Introduction}

Les relations internationales sont fréquemment pensées soit comme un système d'États unitaires
maximisant leur sécurité dans l'anarchie \autocite{waltz1979,mearsheimer2001}, soit comme un
espace de significations et de normes en perpétuelle construction \autocite{wendt1999}. Chacune
de ces lectures capte une part du réel et en manque une autre. Le réalisme structurel saisit la
matérialité de la puissance mais peine à rendre compte du rang, de la légitimité et de la
reconnaissance~; le constructivisme saisit le poids du sens mais sous-estime les capitaux, les
infrastructures et la coercition.

Cet article propose un troisième terme, qui n'est pas un compromis mais une \emph{intégration
disciplinée}~: la Sociologie internationale des personnes collectives (SIPC). Sa thèse centrale
est que la politique mondiale forme une société stratifiée de personnes collectives inégales,
qui mobilisent des capitaux, occupent des rôles, activent des normes, subissent des dépendances
et produisent des trajectoires sous contrainte de légitimité, de crédibilité et de preuve.

L'enjeu n'est pas seulement théorique. La SIPC se veut \emph{opérationnelle}~: elle s'accompagne
d'une méthode d'analyse en dix étapes et d'une ontologie formelle (schémas JSON, validateurs,
exemples), de sorte qu'une analyse puisse être produite, critiquée et \emph{validée par des
règles explicites}. Cette double exigence --- théorique et computationnelle --- distingue la
SIPC d'une simple grille de lecture.

\section{Positionnement~: ce qui est hérité, ce qui est propre}
\label{sec:positionnement}

La SIPC est résolument synthétique. Il importe de dire clairement ce qu'elle doit à d'autres et
ce qu'elle ajoute.

\paragraph{Héritages.} De Bourdieu, elle reprend les notions de champ, de capital et de
conversion, ainsi que l'idée que la position dans un espace structure les stratégies
\autocite{bourdieu1979,bourdieuWacquant1992}. De l'École anglaise, elle retient l'idée d'une
\emph{société} internationale dotée d'institutions \autocite{bull1977,buzan2004}. Du
constructivisme, elle conserve l'attention aux normes, aux identités et à leur dynamique
\autocite{wendt1999,finnemoreSikkink1998,barnettFinnemore2004}. De l'économie politique
internationale, elle hérite l'analyse de l'interdépendance, de la puissance structurelle et de
l'hégémonie \autocite{keohaneNye1977,strange1988,gilpin1987}, prolongée par les travaux récents
sur l'interdépendance instrumentalisée \autocite{farrellNewman2019}. De la sociologie de
l'interaction, elle emprunte la présentation de soi et la \emph{face}
\autocite{goffman1959}, et de la théorie de la reconnaissance la lutte pour le statut
\autocite{honneth1995,ringmar1996}. De la sociologie historique, elle tient l'attention aux
appareils, à la coercition et à la formation des États \autocite{weber1978,tilly1992,mann1986}~;
de Gramsci et Foucault, l'idée que la domination passe par le consentement et la classification
\autocite{gramsci1971,foucault1977}. Enfin, de la sociologie analytique et relationnelle, elle
adopte une explication par \emph{mécanismes} \autocite{hedstromSwedberg1998,elster1989} et une
ontologie processuelle et relationnelle \autocite{emirbayer1997,abbott2016,nexon2009,pouliot2016}.

\paragraph{Apport propre.} La nouveauté de la SIPC ne réside pas dans ces briques mais dans
trois opérations. \emph{Premièrement}, elle les articule autour d'une unité d'analyse unique et
explicite --- la situation sociale internationale --- au lieu de les juxtaposer.
\emph{Deuxièmement}, elle impose une \emph{discipline causale}~: une analyse doit identifier un
mécanisme dominant, un petit nombre de mécanismes secondaires, des trajectoires conditionnelles
et des signaux d'invalidation. \emph{Troisièmement}, elle traduit cet appareil en une
\emph{ontologie formelle} (\S\ref{sec:ontologie}) dont la conformité est vérifiable par des
outils, ce qui transforme une grille conceptuelle en instrument testable.

\section{La personne collective internationale}

Les acteurs internationaux ne sont pas des individus, mais on gagne à les analyser \emph{comme}
des personnes collectives~: des entités composées, institutionnalisées et historiques qui
produisent des décisions par assemblage de bureaucraties, de dirigeants, de doctrines et
d'audiences \autocite{weber1978}. L'analogie doit être maniée avec prudence. Une personne
collective ne possède pas de conscience unitaire~; ce qui lui tient lieu d'« intention » est le
produit instable de coalitions internes et de rapports de force.

L'intérêt analytique de la notion tient à ce qu'elle autorise un vocabulaire commun ---
identité, mémoire, intérêts, ressources, rôles, dépendances, réputation, trajectoire --- sans
homogénéiser les acteurs. Car les personnes collectives sont profondément \emph{inégales}~:
elles diffèrent par leur capacité d'action, leur reconnaissance, leurs capitaux, leur accès aux
institutions, leur crédibilité, leur résilience et, surtout, leur capacité à définir les
catégories du monde \autocite{foucault1977}. La SIPC refuse donc à la fois l'étatisme pauvre
(seuls les États comptent) et le pluralisme naïf (tous les acteurs se valent).

\section{La situation sociale internationale comme unité d'analyse}

L'unité centrale n'est ni l'État, ni la crise, ni le système pris globalement, mais la
\emph{situation sociale internationale}~: une configuration située d'acteurs, de relations, de
champs, de capitaux, de normes, d'institutions, de dépendances, d'audiences et de mécanismes,
qui produit une dynamique observable.

Choisir la crise comme objet rendrait le modèle réactif. Or beaucoup de configurations
décisives demeurent latentes pendant des années~: dépendances énergétiques, rivalités
technologiques, tensions de reconnaissance. La situation permet d'analyser ces configurations
\emph{avant} qu'elles ne deviennent explosives, et de penser la latence comme un état à part
entière et non comme une simple pré-crise.

\section{Champs, capitaux et conversions}

La société internationale n'est pas homogène~: elle se compose de champs --- militaire,
diplomatique, économique, financier, technologique, juridique, institutionnel, symbolique,
informationnel --- dotés chacun de règles, de ressources et de formes de reconnaissance propres
\autocite{bourdieu1979}. Un acteur peut être puissant dans un champ et faible dans un autre.

La notion décisive est la \emph{conversion de capital}. Un capital n'est stratégique que s'il
peut être converti en effet dans une situation donnée~: le capital militaire en dissuasion ou
coercition, le capital financier en sanctions ou en discipline \autocite{strange1988}, le
capital technologique en dépendance ou en verrouillage, le capital symbolique en légitimité ou
en mobilisation. D'où une définition relationnelle de la puissance~:

\begin{quote}
La puissance n'est pas la simple possession de ressources~; elle est la capacité \emph{située}
de convertir des capitaux en effets reconnus.
\end{quote}

\section{Rôles, rangs et audiences}

Les personnes collectives n'agissent pas seulement~: elles se présentent, se justifient et sont
classées par d'autres \autocite{goffman1959}. Un acteur peut revendiquer un rôle (protecteur,
puissance responsable, gardien de norme, victime) sans qu'il soit reconnu. La SIPC distingue
ainsi rôle revendiqué, reconnu, contesté, imposé, performé et discrédité. La même logique vaut
pour le rang~: une part majeure de la politique internationale est une lutte pour la
reconnaissance du statut \autocite{honneth1995,ringmar1996,pouliot2016}. Les classements
internationaux (« puissance responsable », « État voyou », « marché émergent ») ne décrivent pas
seulement~; ils produisent des effets d'accès, d'exclusion et de réputation, et constituent
eux-mêmes des instruments de pouvoir.

\section{Normes, institutions et légitimité}

Une norme internationale est une attente stabilisée définissant l'acceptable, l'obligatoire et
l'interdit \autocite{finnemoreSikkink1998}. Mais une norme n'agit pas seule~: elle doit être
\emph{activée} par un acteur qui l'invoque sur une scène, devant une audience qui reconnaît ou
conteste, avec un coût ou une sanction à la clé. Une norme peut protéger comme hiérarchiser,
être sincère et instrumentalisée, universelle dans son langage et sélective dans son
application.

Les institutions, quant à elles, sont des dispositifs sociaux de stabilisation, de classement
et de lutte \autocite{barnettFinnemore2004}. Elles cristallisent souvent un rapport de force en
procédure~: une institution naît d'un rapport de force, le transforme en règle, et la règle peut
ensuite survivre au rapport de force initial. La légitimité n'est donc jamais une propriété
fixe~: c'est une reconnaissance située, disputée et convertible.

\section{Dépendances, vulnérabilités et résilience}

La SIPC accorde une place centrale aux dépendances. Une dépendance est une relation où un acteur
a besoin d'un autre, d'un réseau ou d'une infrastructure pour maintenir une capacité d'action.
Contre deux illusions symétriques --- « dépendance = faiblesse » et « interdépendance = paix »
\autocite{keohaneNye1977} --- la SIPC soutient qu'une interdépendance peut tout autant stabiliser
que devenir un levier de coercition \autocite{farrellNewman2019}. Une dépendance se qualifie par
sa criticité, son asymétrie, sa substituabilité, son coût de rupture et le potentiel de sa
\emph{weaponisation}. La résilience est alors la capacité à absorber et reconfigurer ses
dépendances sans perdre ses fonctions essentielles~: l'autonomie stratégique n'est pas l'absence
de dépendance, mais la maîtrise politique des dépendances vitales.

\section{Domination, contestation et ordre}

L'hégémonie tient rarement par la seule coercition. Elle tient lorsqu'une domination matérielle
est convertie en institutions, normes, dépendances, protections, classements et consentement
partiel \autocite{gramsci1971,gilpin1987}. Signe révélateur de stress~: quand une hégémonie doit
menacer de plus en plus pour obtenir ce qu'elle obtenait par simple normalité, elle est déjà en
difficulté. La contestation, de son côté, vise à modifier règles, rangs, dépendances ou
catégories~; mais la SIPC pose une règle froide~: la capacité de nuisance n'est pas une capacité
d'ordre. Un acteur peut affaiblir un ordre sans pouvoir en produire un autre. Un ordre, enfin,
n'est ni la paix ni la justice~: c'est la prévisibilité relative des positions, des règles, des
coûts et des accès. Son changement est le plus souvent une recomposition multi-couches, non un
effondrement brutal.

\section{Causalité configurationnelle et mécanismes}
\label{sec:causalite}

La SIPC n'explique pas par une cause unique mais par une \emph{configuration}~: un effet résulte
d'une configuration initiale, de mécanismes activés, de seuils franchis, d'une séquence
d'interactions, de rétroactions et de verrouillages \autocite{hedstromSwedberg1998,elster1989}.
Un mécanisme est une chaîne causale activable sous conditions, dotée d'une séquence typique et
de signaux d'invalidation. La SIPC en retient une dizaine~: conversion de capital, weaponisation
de dépendance, verrouillage d'audience, crise de reconnaissance, dilemme de sécurité
\autocite{jervis1976}, activation normative, lutte de classement, écart de crédibilité,
prophétie auto-réalisatrice, désynchronisation d'ordre.

La règle méthodologique est stricte~: pas de trajectoire sans mécanisme, pas de mécanisme sans
conditions d'activation, pas d'analyse sérieuse sans signal d'invalidation. Une analyse
identifie un mécanisme \emph{dominant}, deux à quatre mécanismes secondaires au plus, des seuils,
une séquence et les signaux qui feraient changer le diagnostic.

\section{Preuve, contre-hypothèses et invalidation}
\label{sec:preuve}

Le danger majeur d'un tel cadre est de produire des interprétations brillantes mais
insuffisamment prouvées. La parade est la \emph{triangulation probatoire}
\autocite{beachPedersen2013}~: une hypothèse n'est robuste que si elle relie sources
déclaratives, comportements observables, données matérielles, séquence temporelle, réactions
d'audience, contre-hypothèses et contradictions, le tout assorti d'un niveau de confiance
explicite. La SIPC recommande une échelle de statuts probatoires (de \textsc{speculative} à
\textsc{strongly\_supported}) et impose qu'aucun score ne vaille sans justification textuelle.

\begin{quote}
Un diagnostic incertain n'est pas un échec~; un diagnostic faussement certain est une faute.
\end{quote}

La SIPC renforce cette discipline par quatre dispositifs. \emph{Premièrement}, une
\textbf{calibration probabiliste ancrée}~: les bandes verbales (de \textsc{low} à \textsc{high})
sont reliées à des plages numériques fixes, et toute estimation chiffrée doit y être cohérente,
ce qui prévient la fausse précision. \emph{Deuxièmement}, une \textbf{analyse des hypothèses
concurrentes} (ACH) à la Heuer~: on construit une matrice hypothèses~$\times$~preuves et l'on
retient l'hypothèse la \emph{moins infirmée} plutôt que la plus confirmée, en pondérant chaque
preuve par sa \emph{diagnosticité} et son type de test de \emph{process tracing} (indice,
\emph{hoop}, \emph{smoking gun}, doublement décisif) \autocite{beachPedersen2013}. On y adjoint
un \emph{Key Assumptions Check} et un \emph{premortem} comme garde-fous de débiaisage.
\emph{Troisièmement}, une \textbf{traçabilité des sources} qui sépare la fiabilité de la source
(échelle A--F) de la crédibilité de l'information (1--6), selon le code Admiralty, afin que
« source » ne soit jamais confondu avec « preuve ». \emph{Quatrièmement}, une \textbf{boucle de
rétroaction}~: les signaux d'invalidation sont surveillés, les analyses versionnées, et le
dénouement des trajectoires fait l'objet d'un \emph{backtesting} par score de Brier, seule voie
vers une calibration empirique progressive. Ces exigences sont graduées selon la maturité de
l'analyse et vérifiées automatiquement.

\section{Études de cas~: le détroit de Taïwan et un corpus de cinq situations}
\label{sec:taiwan}

Nous appliquons l'appareil à une situation qui active presque toutes ses briques.

\paragraph{Cadrage.} La situation du détroit de Taïwan n'est pas réductible à une crise
militaire potentielle. C'est une situation sociale internationale durable, périodiquement
accélérée, qui articule souveraineté contestée, reconnaissance partielle, dépendance mondiale
aux semi-conducteurs avancés, dissuasion sino-américaine et ordre régional sous stress. Les
champs activés sont au premier chef militaire, technologique, diplomatique et économique.

\paragraph{Acteurs et relations.} On retient cinq personnes collectives principales~: la Chine
(revendication de souveraineté, capitaux militaire, économique, technologique et symbolique)~;
les États-Unis (garant d'un ordre régional sous ambiguïté stratégique)~; Taïwan (entité de
\emph{facto} autonome, statut contesté, nœud technologique)~; TSMC (firme détenant un quasi-
monopole sur les nœuds de gravure les plus avancés)~; le Japon (allié exposé). Les relations
structurantes sont une non-reconnaissance coercitive (Chine~$\rightarrow$~Taïwan), une protection
sous ambiguïté (États-Unis~$\rightarrow$~Taïwan), une rivalité stratégique (États-Unis~$\leftrightarrow$~Chine)
et une dépendance technologique mondiale au nœud taïwanais.

\paragraph{Capitaux et dépendances.} Le capital décisif de Taïwan n'est pas militaire mais
technologique~: la concentration de la production de pointe constitue un \emph{chokepoint}
mondial. Cette dépendance est doublement asymétrique et faiblement substituable à court terme~;
son potentiel de weaponisation est réel mais borné par le coût d'une rupture pour tous les
acteurs, y compris la Chine, elle-même dépendante d'intrants technologiques occidentaux.

\paragraph{Mécanisme dominant et mécanismes secondaires.} Le mécanisme dominant est une
\emph{crise de reconnaissance}~: l'écart entre le statut revendiqué et la reconnaissance
accordée structure la tension et rend certaines concessions coûteuses devant les audiences
internes \autocite{ringmar1996}. Deux mécanismes secondaires l'amplifient~: un \emph{dilemme de
sécurité} sino-américain, où chaque mesure défensive est lue comme offensive
\autocite{jervis1976}, et une \emph{weaponisation potentielle de la dépendance} technologique
\autocite{farrellNewman2019}. Les amortisseurs --- coût d'une guerre majeure, interdépendance
économique, canaux de gestion de crise --- maintiennent ordinairement la situation sous le seuil
de rupture.

\paragraph{Trajectoires.} Deux trajectoires conditionnelles se dégagent. La première, jugée la
plus probable, est une \emph{tension gérée} sans rupture majeure, conditionnée par la continuité
de la dissuasion, l'absence d'incident accidentel grave et la reconnaissance du coût économique.
La seconde est une \emph{escalade coercitive} (blocus ou quarantaine), conditionnée par un
verrouillage d'audience et un incident déclencheur. Les points de bifurcation incluent un
incident militaire majeur, un basculement politique interne et une expansion des contrôles
technologiques.

\paragraph{Preuve et contre-hypothèses.} L'hypothèse du primat de la reconnaissance s'appuie sur
un faisceau convergent~: réaffirmations officielles récurrentes (discours), exercices et
incursions répétés (pratique observable), concentration documentée de la production de
semi-conducteurs (donnée matérielle). Elle doit être confrontée à au moins deux
contre-hypothèses~: que le moteur premier soit le dilemme de sécurité, la reconnaissance n'étant
qu'un habillage~; ou que la dépendance technologique soit la variable décisive. Une matrice ACH
départage ces hypothèses~: une preuve à forte diagnosticité (le traitement de l'indépendance
formelle comme \emph{casus belli} indépendamment des arrangements de sécurité, de type
\emph{smoking gun}) infirme à la fois l'hypothèse de pur dilemme de sécurité et celle de pure
dépendance, laissant la reconnaissance comme hypothèse la moins infirmée. Un signal
d'invalidation clair serait la résolution durable de la tension par des arrangements de sécurité
\emph{sans} composante de statut. Le statut probatoire raisonnable demeure
\textsc{moderately\_supported}, avec une confiance moyenne~: un diagnostic prudent, conforme à la
discipline de la \S\ref{sec:preuve}.

Ce cas illustre la valeur ajoutée de la SIPC~: il ne se contente pas de lister des facteurs, il
les \emph{hiérarchise} autour d'un mécanisme dominant, les ordonne en trajectoires
conditionnelles, et expose les conditions auxquelles l'analyse devrait être révisée.

\paragraph{Un corpus de cas validés.} Au-delà de Taïwan, l'appareil a été appliqué à cinq
situations contrastées, toutes formalisées comme objets JSON conformes au tier le plus exigeant
(\textsc{validated}), chacune mettant en avant un mécanisme dominant différent~:

\begin{itemize}[nosep]
  \item \textbf{détroit de Taïwan} --- \emph{crise de reconnaissance} (statut contesté,
  dissuasion, dépendance aux semi-conducteurs)~;
  \item \textbf{guerre en Ukraine} --- \emph{activation normative} autour de la souveraineté,
  la résilience comparée des coalitions devenant la variable critique~;
  \item \textbf{rivalité technologique États-Unis / Chine} --- \emph{weaponisation de
  chokepoints} (lithographie avancée) produisant un découplage sélectif~;
  \item \textbf{mer de Chine méridionale} --- \emph{conversion incrémentale} d'un capital
  paramilitaire en contrôle de fait, calibrée sous le seuil d'activation d'une alliance~;
  \item \textbf{Arctique} --- \emph{désynchronisation d'ordre}, l'ouverture physique et
  économique dépassant l'adaptation des institutions et créant un déficit de gouvernance.
\end{itemize}

Cette variété est un test de portée~: le même appareil discrimine des moteurs causaux distincts
sans se réduire à une grille passe-partout, et chaque analyse expose une matrice d'hypothèses
concurrentes où l'hypothèse retenue est la moins infirmée.

\section{Formalisation~: une ontologie testable}
\label{sec:ontologie}

L'appareil se traduit en une ontologie JSON dont l'objet racine, \texttt{situation\_analysis},
agrège des profils d'acteurs, de relations, de champs, de capitaux, de dépendances, de normes,
d'institutions, de mécanismes, de trajectoires, de preuves et de diagnostic. Deux niveaux de
validation accompagnent l'ontologie~: une validation \emph{structurelle} par schémas JSON, et
une validation \emph{doctrinale} qui encode les règles non exprimables par le schéma --- présence
d'un mécanisme dominant, d'au moins une contre-hypothèse, de signaux d'invalidation, intégrité
référentielle entre objets, et interdiction des scores non justifiés. Cette validation
doctrinale est \emph{graduée} selon la maturité de l'analyse (\textsc{draft},
\textsc{review\_required}, \textsc{validated})~: les tiers supérieurs exigent une ACH cohérente,
des sources gradées (profil \texttt{source\_profile}) et un suivi des signaux d'invalidation.
Les cinq études de cas de la \S\ref{sec:taiwan} existent ainsi sous forme d'objets JSON
validés au tier le plus exigeant, accompagnés d'outils qui vérifient automatiquement leur
conformité (validation des schémas, contrôle doctrinal, rendu de la matrice ACH, backtesting).
C'est cette vérifiabilité qui distingue la SIPC d'une grille purement discursive.

\section{Limites et risques}

La SIPC affiche ses limites, condition de son sérieux. Elle ne produit pas de certitude et ne
prédit pas mécaniquement l'avenir~; elle ne remplace ni le droit international, ni l'économie
politique, ni la stratégie. Trois risques internes la guettent. Le premier est
l'\emph{anthropomorphisme}~: prêter une intention unitaire à un composite~; on le contient en
rappelant que les attributs de la personne collective sont produits par des appareils. Le
deuxième est l'\emph{hypertrophie conceptuelle}~: la multiplication des notions sans hiérarchie~;
on le contient par la discipline du mécanisme dominant. Le troisième est la \emph{fausse
précision}~: l'apparence de rigueur des scores numériques non calibrés~; on le contient en
maintenant les scores prudents et toujours justifiés. Enfin, le cadre n'échappe pas à la
\emph{réflexivité}~: les catégories qu'il mobilise sont elles-mêmes des instruments de
classement, ce que l'analyste doit assumer.

\section{Conclusion}

La SIPC propose de lire la politique mondiale comme une société de personnes collectives
inégales, et d'organiser cette lecture autour d'une unité d'analyse claire, d'une causalité par
mécanismes et d'une discipline probatoire. Sa singularité n'est pas conceptuelle mais
architecturale~: intégrer des héritages dispersés en un appareil cohérent, l'astreindre à une
calibration probabiliste ancrée, à une analyse des hypothèses concurrentes, à une traçabilité
des sources et à une boucle de suivi et de backtesting, puis le rendre \emph{testable} par des
outils de conformité gradués. Les cinq études de cas --- détroit de Taïwan, guerre en Ukraine,
rivalité technologique États-Unis / Chine, mer de Chine méridionale et Arctique --- suggèrent
qu'un tel cadre peut hiérarchiser les causes, départager des hypothèses concurrentes, expliciter
les trajectoires et nommer ses propres conditions d'invalidation --- ce qui, pour une sociologie
appliquée des relations internationales, constitue déjà une exigence salutaire.

\printbibliography[heading=bibintoc,title={Bibliographie}]

\end{document}
